\documentclass[11pt]{article}
\usepackage{amsmath,amssymb,amsthm}
\usepackage[margin=1in]{geometry}
\usepackage[hidelinks]{hyperref}
\newtheorem{theorem}{Theorem}
\newtheorem{proposition}[theorem]{Proposition}
\newtheorem{corollary}[theorem]{Corollary}
\newcommand{\C}{\mathcal C}
\newcommand{\M}{\mathfrak M}
\newcommand{\MA}{\operatorname{MA}}
\newcommand{\V}{\mathcal V}
\newcommand{\R}{\mathbb R}
\setlength{\parindent}{0pt}
\setlength{\parskip}{0.5em}
\emergencystretch=3em
\title{Affine dual rays and quantitative separation of variational measure limits}
\author{Unreviewed scoped deductions for original P16.7}
\date{9 October 2026}
\begin{document}
\maketitle
\begin{abstract}
Within the explicitly qualified bounded calculus on the actual analytic
torus quotient, the full dual is affine along the segment from any
eligible prescribed measure to its canonical variational weak limit.
Even approximate maximizers with a varying segment parameter have the
same energy, full-point values and weak Monge--Ampere limit.
For two arbitrary eligible measures, equality in dual convexity is
equivalent to a common maximizing family. Distinct variational measure
limits force a quantitatively positive Jensen gap, detected by actual
continuous positive weights. These are consequences of the existing
qualified calculus, not proofs of a new analytic representative or of
the prescribed-measure equation. Original P16.7 remains open.
\end{abstract}

\section{Actual domain and exact qualified inputs}
Use the actual compact Hausdorff analytic surface
\[
 K=\mathbb C((s)),\quad |s|=e^{-1/8},\quad t=s^8,\qquad
 X=(\mathbb G_m^2)^{\mathrm{an}}_K/\langle(t,2),(3,t)\rangle.
\]
The continuous class $\C$ is the original minimal continuous CLD
positive-weight class with local reference $\rho=|\tau|^2/2$,
$\theta=d'd''\rho$, raw mass $\int\theta^2=2$, and actual bounded
products $\MA(a)=(\theta+d'd''a)^2$. The energy is
\[
 E(a)=\frac13\int_Xa\,[\theta^2+
       \theta\wedge(\theta+d'd''a)+\MA(a)].
\]
Define $F_\lambda(a)=E(a)-\int a\,d\lambda$ and the full bounded
dual $B(\lambda)=\sup_{a\in\C}F_\lambda(a)$. Let $\M$ be all
positive raw-mass-two Radon measures with $B(\lambda)<\infty$,
with no support or domination restriction. Constants leave $F$ unchanged;
all maximizing weights below are normalized by $\sup_Xa=0$.

The exact inputs are the complete, separately qualified proofs in
ROOT414 and ROOT415. They have no automatic correctness grant.
For $b\ge0$ put
\[
 C(b)=\tfrac14(\sqrt{6b}+\sqrt{10b+8})^2,\qquad
 \C_C=\{a\in\C:\sup a=0,\ E(a)\ge-C\}.
\]
If $B(\lambda)\le b$ and $F_\lambda(a)\ge-1$, then $a\in\C_{C(b)}$.
For $a,c\in\C_C$, the nonnegative symmetric quantity
\[
 I(a,c)=\int_X(a-c)(\MA(c)-\MA(a))
\]
has exact midpoint energy gain $I(a,c)/8$. Consequently, with
$\delta_\lambda(a)=B(\lambda)-F_\lambda(a)\ge0$,
\begin{equation}\label{mid}
 I(a,c)\le4(\delta_\lambda(a)+\delta_\lambda(c)).
\end{equation}
This inequality needs no maximizer attainment and holds for any two
bounded competitors, by applying the objective bound to their midpoint.
If $h\in\C$, subtract its supremum and set
$L=\max(C,-E(h-\sup h),0)$. The actual Green estimate is
\begin{equation}\label{green}
 \left|\int_Xh\,[\MA(a)-\MA(c)]\right|
       \le\sqrt{12L\,I(a,c)}.
\end{equation}
Indeed set $T=\theta_a+\theta_c$ and $f=a-c$. Product
transgression and Green identify the integral as $-Q_T(h,f)$,
$Q_T(f,f)=I(a,c)$, and the bounded gradient estimate gives
$Q_T(h,h)\le12L$. Actual Cauchy--Schwarz proves~\eqref{green}.
The span of $\C$ is uniformly dense in $C(X,\R)$ by the actual
Laurent-cap separation and compact vector-lattice proof.

On each $\C_C$, $I(a,c)\to0$ controls differences of real energies,
reference and self mixed moments, diagonal MA moments, each fixed
bounded-test $L^1$ distance, and each full graded-dimension-two point
value. Write $H_C$ for the real-energy modulus and
$P_{C,x}$ for the full-point modulus; both tend to zero at zero.
The diagonal modulus is $H_C(q)+\sqrt{12Cq}$.
ROOT415 also gives a uniform third-slot absolute moment modulus
$T_C(q)\to0$ for every $w\in\C_C$.

Every normalized maximizing net for $\lambda\in\M$ has a canonical
weak MA limit $\V\lambda\in\M$, independent of the family. The same
net maximizes $B(\V\lambda)$; its limiting real energy and diagonal
moment $e,p$ satisfy $B(\V\lambda)=e-p$. Furthermore
\begin{equation}\label{retr}
 B(\V\lambda)\le B(\lambda),\qquad \V(\V\lambda)=\V\lambda.
\end{equation}
ROOT415 constructs the canonical ordinary USC upper relaxation from
all affinoid Shilov limits. It is determined by the full-point tuple;
its positivity, original finite-energy membership and MA equation
remain unresolved. We use its canonical formula only as a topological
invariant, not as an analytic solution.

These inputs were read at the current applications: the whole467-line
ROOT414 proof and ROOT415 universal-limit, dual-calibration, defect and
USC-relaxation proofs. Their inherited CLD/GR, generated, geometric,
nonalgebraization and deeper source qualifications remain unchanged.
This is not a fresh recursive primary audit. Compiler records and
literal identity do not certify mathematical correctness.

\section{The entire segment to the variational limit}
\begin{theorem}\label{ray}
Let $\mu\in\M$, $\nu=\V\mu$, $B=B(\mu)$, $b=B(\nu)$, and
$\mu_t=(1-t)\mu+t\nu$, $0\le t\le1$. Then
\[
 B(\mu_t)=(1-t)B+tb,\qquad \V\mu_t=\nu.
\]
All these measures have the same canonical real energy, reference and
self mixed moments, diagonal moment, full-point tuple and ordinary USC
upper relaxation. Their dual defect is exactly
$B(\mu_t)-B(\V\mu_t)=(1-t)(B-b)$.

More strongly, let $t_i\in[0,1]$ vary arbitrarily along a directed net
and let normalized $c_i\in\C$ satisfy
$\epsilon_i=B(\mu_{t_i})-F_{\mu_{t_i}}(c_i)\to0$.
Then $\MA(c_i)\rightharpoonup\nu$ against every continuous test,
and the listed real moments and each full-point value converge to
their common canonical limits, without requiring a limit of $t_i$.
\end{theorem}
\begin{proof}
Choose a normalized maximizing net $a_j$ for $\mu$. By exact
calibration it also maximizes $\nu$. Convexity gives
$B(\mu_t)\le(1-t)B+tb$; evaluating on $a_j$ and passing to its
two objective limits gives the reverse inequality. Thus equality holds
for every $t$, including the endpoints, and $a_j$ maximizes each
$\mu_t$. Its actual MA measures already tend to $\nu$, so canonical
family independence gives $\V\mu_t=\nu$. Applying the same family
to all the established canonical invariants proves their equality.
The displayed defect follows by subtraction.

For the varying-parameter claim set
$\eta_j=\max(\delta_\mu(a_j),\delta_\nu(a_j))\to0$. Exact dual
affinity gives
\[
 (1-t_i)\delta_\mu(c_i)+t_i\delta_\nu(c_i)=\epsilon_i.
\]
At each index choose $\lambda_i=\mu$ if $t_i\le1/2$, and
$\lambda_i=\nu$ otherwise. Nonnegative deficits imply
$\delta_{\lambda_i}(c_i)\le2\epsilon_i$. This choice is used only
in the estimate; no cofinality of either subset of indices is presumed.
All $B(\mu_{t_i})\le B$, so universal coercivity gives a common
late $\C_{C(B)}$ for $c_i$ and $a_j$. Equation~\eqref{mid} gives
\begin{equation}\label{uniformray}
 I(c_i,a_j)\le8\epsilon_i+4\eta_j.
\end{equation}
For fixed $h\in\C$ pass $j$ to infinity in~\eqref{green}, using
the genuine weak convergence of $\MA(a_j)$. It follows that
\[
 \left|\int h\,[\MA(c_i)-\nu]\right|
       \le\sqrt{96L\epsilon_i}\longrightarrow0.
\]
Uniform density of the span and common mass two extend convergence
to every continuous test; the two replacement errors total at most
$4\|h-h_0\|_\infty$. Equation~\eqref{uniformray} and the real,
diagonal and mixed-moment moduli identify each listed scalar limit
by first choosing $j$ late, then $i$ late. The same two-stage argument
with $P_{C,x}$ proves convergence at each full point $x$.
No pointwise limit at other points or positivity of an upper
regularization is inferred.
\end{proof}

\section{Affine dual segments and a measured Jensen gap}
\begin{proposition}\label{affine}
For $\mu_0,\mu_1\in\M$, equality
\[
 B((1-t)\mu_0+t\mu_1)=(1-t)B(\mu_0)+tB(\mu_1)
\]
at any one $0<t<1$ is equivalent to the existence of a common
normalized maximizing sequence for both endpoint measures.
It implies equality for every $t\in[0,1]$, equality of their
variational measure limits and all canonical invariants above.
The converse from equality of variational limits alone is not asserted.
\end{proposition}
\begin{proof}
If equality holds at $t$, take a normalized maximizing sequence $c_n$
for the mixture; such a sequence exists from its finite real supremum.
The exact deficit identity reads
\[
 (1-t)\delta_{\mu_0}(c_n)+t\delta_{\mu_1}(c_n)
 =B((1-t)\mu_0+t\mu_1)-F_{(1-t)\mu_0+t\mu_1}(c_n)\to0.
\]
Both summands are nonnegative with fixed positive coefficients;
each endpoint deficit tends to zero. Thus the sequence maximizes
both endpoints. Conversely a common maximizing sequence gives the
convexity reverse inequality at every parameter. Its unique actual
weak MA limit and all canonical limits agree for both endpoints.
\end{proof}

\begin{theorem}\label{jensen}
Put $\nu_k=\V\mu_k$, $B_k=B(\mu_k)$,
$\mu_t=(1-t)\mu_0+t\mu_1$, and $0<t<1$. Define
\[
 J_t=(1-t)B_0+tB_1-B(\mu_t)\ge0,\qquad
 C=C(\max(B_0,B_1)).
\]
For each actual $h\in\C$ set $L=\max(C,-E(h-\sup h),0)>0$.
Then
\begin{equation}\label{lower}
 J_t\ge\frac{t(1-t)}{48L}
       \left|\int_Xh\,(\nu_0-\nu_1)\right|^2.
\end{equation}
In particular, $\nu_0\ne\nu_1$ forces strict dual convexity
at every interior parameter, detected by at least one $h\in\C$.
\end{theorem}
\begin{proof}
Let $c_i$ maximize $\mu_t$ and let $a_j^{(k)}$ maximize $\mu_k$.
Normalize all three nets. The mixture has finite dual bounded by
$\max(B_0,B_1)$, so their late tails share $\C_C$. Write
$\epsilon_i=B(\mu_t)-F_{\mu_t}(c_i)\to0$. The exact identity is
\[
 (1-t)\delta_{\mu_0}(c_i)+t\delta_{\mu_1}(c_i)=J_t+\epsilon_i.
\]
Therefore $\delta_{\mu_0}(c_i)\le(J_t+\epsilon_i)/(1-t)$ and
$\delta_{\mu_1}(c_i)\le(J_t+\epsilon_i)/t$. For each fixed $i$,
apply~\eqref{mid} at endpoint $k$, then~\eqref{green}, and pass
$j$ to infinity using weak convergence against the fixed continuous
$h$. This yields
\[
 \left|\int h\,[\MA(c_i)-\nu_k]\right|
       \le\sqrt{48L\delta_{\mu_k}(c_i)}.
\]
The triangle inequality at this fixed index, followed by weighted
Cauchy--Schwarz, gives
\[
 \left|\int h\,(\nu_0-\nu_1)\right|
 \le\sqrt{48L}\bigl(\sqrt{\delta_{\mu_0}(c_i)}+
                       \sqrt{\delta_{\mu_1}(c_i)}\bigr)
 \le\sqrt{\frac{48L(J_t+\epsilon_i)}{t(1-t)}}.
\]
Indeed the weighted squares sum to $J_t+\epsilon_i$, while the
reciprocal weights sum to $(1-t)^{-1}+t^{-1}=1/(t(1-t))$.
Pass $i$ to infinity and square to obtain~\eqref{lower}. This
part needs no convergence of $\MA(c_i)$.
Here $C(b)\ge C(0)=2$, so the denominator is strictly positive.
If the two actual measures differ, some continuous test distinguishes
them. Density of the span implies that some member of $\C$ does
so: otherwise all span integrals, and hence all continuous integrals,
would agree. Its lower bound is positive at every $0<t<1$.
\end{proof}

\begin{corollary}
The full dual is strictly convex along every segment with distinct
endpoints that are already actual generated MA measures. Conversely,
if the full dual is strictly convex on the entire eligible measure
class $\M$, then $\V\mu=\mu$ for every $\mu\in\M$.
Neither statement supplies analytic representation of an arbitrary
fixed point.
\end{corollary}
\begin{proof}
ROOT414 proves that every actual generated MA measure is fixed by
$\V$. Thus distinct such endpoints have distinct images, and
Theorem~\ref{jensen} applies. For the converse,
Theorem~\ref{ray} makes $B$ affine between $\mu$ and $\V\mu$;
strict convexity on $\M$ forces those endpoints to coincide.
\end{proof}

\section{Exact scope and remaining original requirements}
Theorem~\ref{ray} makes each fibre of $\V$ star-shaped about its
fixed-point measure, and gives uniform approximate-maximizer control
on those specific rays. It does not prove convexity of a whole fibre,
or that equal images imply dual affinity. The Jensen lower bound is
a genuine consequence of the actual Green estimate and bounded
maximizing calculus; no abstract inverse MA operator is introduced.

Neither a zero defect, an affine ray, an idempotent measure map nor
the common USC relaxation constructs an actual positive analytic
weight in the original finite-energy domain. Equality $\V\mu=\mu$
for every eligible measure, recognition of the canonical upper
relaxation, sufficient envelopes or legitimate two-sided variations,
and the full prescribed-measure existence and uniqueness remain open.
All inherited analytic/geometric/source qualifications are preserved.
This paper is UNREVIEWED; it closes no original atom and makes no
new scoped acceptance. Counts remain16/42 original problems solved,
26 open and97 historically accepted scoped papers. Primary literal
bodies, pixels and books stay LOCALONLY. The protected P16.4 editor
remains open and unchanged; active native289/290 mathematics is not
read or hot-fed by this independent deduction.
\end{document}
